\documentclass[11pt]{article}
\usepackage{amsmath,amssymb,amsthm}
\usepackage[margin=1in]{geometry}
\usepackage{hyperref}
\newtheorem{theorem}{Theorem}
\newtheorem{lemma}[theorem]{Lemma}
\newtheorem{proposition}[theorem]{Proposition}
\newcommand{\R}{\mathbb R}
\newcommand{\Z}{\mathbb Z}
\newcommand{\Int}{\operatorname{int}}
\newcommand{\Gr}{\operatorname{Gr}}
\newcommand{\T}{\mathcal T}
\title{The full one-critical-point obstruction over \(S^9\)}
\author{Campaign 9000340: Engineer argument for independent review}
\date{November 5, 2026}
\begin{document}
\maketitle

\paragraph{Result and evidence boundary.}
We verify Proposition \texttt{one-critical} of the staged September 23,
2026 manuscript, under its actual varying-infinity hypotheses.
The new work supplies the negative-eigenbundle comparison, quantitative
normalization and smoothing, and a collar-relative, parameter-uniform
straightening of \emph{unframed} tube cores.
The previously independently accepted result supplies only the normalized
zero-band deduction; it is stated precisely below and is not reproved.
This is a prose proof using identified external theorems, supported by an
exact calculation certificate. It is not a kernel-checked topological
proof, a counterexample, or a novelty certification.
The nonzero tube-kernel calculation and the manuscript's Lagrangian
realization and transport remain unverified by this note.

\section{Statement, conventions, and inputs}

Write \(q_{k,l}(x,y)=-|x|^2+|y|^2\).
Let \(T_{k,l}\) be the \(C^1\) path component of
\(q_{k,l}|_{S^{k+l-1}}\) among functions having zero as a regular value.
Its elements extend homogeneously by \(g(ru)=r^2h(u)\), including their
\(C^1\) extension at zero. Smooth homogeneous data mean that \(h\) is
smooth. The stable tube space \(\T\) uses the diagonal stabilization
by \(q_{1,1}\). The spherical class of a family is that of its
negative domains \(D(h)=\{h\leq0\}\), stabilized as sphere homotopy types.
This definition does not require the tube-fibration theorem.

\begin{theorem}\label{thm:main}
Let \(p\geq20\), \(B=S^9\), and let
\(f:B\times\R^{2p}\longrightarrow\R\) be smooth.
Suppose that outside one common ball it agrees with smooth
degree-two homogeneous data \(g_b\) of type \(T_{p,p}\);
that each \(f_b\) has exactly one nondegenerate critical point,
forming a smooth section; and that its stable spherical infinity class
is zero in \([B,BG]\).
Then its stable tube infinity class is zero in \([B,\T]\).
\end{theorem}

Here are the imported results and their exact scope.
\begin{enumerate}
\item \textbf{Accepted zero-band result.}
For a smooth family over \(S^9\), with just the critical point zero,
equal to the \emph{fixed} \(q_{k,l}\) on a common neighborhood of zero,
and smooth homogeneous \(T_{k,l}\) data at infinity, \(k,l\geq20\),
the negative flow embeds the fixed inner truncated tube \(T^0\) into
the outer truncated tube \(T_b^1\). Its complement is an
incoming-identified h-cobordism family \(C_b\), and this family is a
product relative to its incoming collar.
This is the result of \path{research/zero_band_audit/proof.tex},
especially its h-cobordism lemma and vanishing theorem.
The Host-invoked Reviewer accepted precisely this deduction in
\path{handoffs/5e635ae1a6ab/round-0001.json} of the assigned project.
The seal does not approve any of the new steps in this note.
\item \textbf{Stable \(J_8\) is injective.}
Adams \cite[Theorem 9.5, p.~51]{Adams} states that, for \(r>1\) and
\(r\equiv0,1\pmod8\),
\[
 e_RJ:\pi_r(SO)\longrightarrow\Z/2
 \quad\text{is an isomorphism}.
 \]
Here \(J\) is the stable map defined on p.~21 and \(e_R\) is the
mod-two invariant of Section 9, not the differently graded
\(\mathbb Q/\mathbb Z\)-valued invariant. In degree eight its domain
is all of \(\pi_8^S\), since the real \(K\)-theory invariant \(d_R\)
vanishes there (Section 7).
Specializing the displayed composition to \(r=8\) proves injectivity
of \(J_8:\pi_8(O)\to\pi_8^S\); positive-degree homotopy is unchanged
on passing from \(SO\) to \(O\).
The map on \(\pi_9\) induced by \(BO\to BG\) is this \(J_8\).
Neither \(J_9\), Rognes's computation, nor the existence of a nonzero
kernel class in \(\pi_9\T\) is an input.
\item \textbf{Homological Whitehead.}
An integral homology isomorphism between simply connected CW
complexes is a homotopy equivalence
\cite[Corollary 4.33]{Hatcher}. The retained author text and the
independent source reading from the zero-band review are reused.
\item \textbf{Multijet transversality.}
For a smooth submanifold \(S\) of \(J_s^r(X,Y)\), smooth maps whose
multijet extension is transverse to \(S\) form a residual subset
\cite[Theorem 2.3.4]{BPT}, in the smooth topology.
The cited version's Corollary 2.3.5(iii) also gives the
\emph{non-relative} fiberwise parameter count
\(\operatorname{codim}S>s\dim X+d\) for disk parameters of
dimension \(d\).
It does not prescribe endpoint values or collars.
Lemma~\ref{lem:relative} below derives the relative assertion we
actually need, from the ordinary multijet theorem on closed manifolds.
\end{enumerate}

The Adams and BPT original PDFs and searchable texts are retained in
\texttt{sources/}, with the exact source identifiers in the claim ledger.
The BPT file is explicitly arXiv version 1.
Bott's original paper was also obtained: its Corollary to Theorem II,
equation (1.5), and the following orthogonal period give
\(\pi_8(O)=\Z/2\). That independent group confirmation is not needed
in addition to Adams's stated isomorphism, and is not a hidden
kernel-class assumption.
The formal citation-archive operation writes outside this mission's
owned paths. It has not been used or represented as complete.

We use elementary smooth ODE theory, the inverse function theorem,
smooth partitions of unity, basic CW obstruction theory and homology.
In particular, none of the arguments below select a family of
trivializations from merely pointwise trivializations.

\section{Compact normalization and the negative sphere bundle}

For smooth angular data, direct differentiation gives
\[
 \nabla g_b(ru)=r\bigl(2h_b(u)u+\nabla_Sh_b(u)\bigr),\qquad
 |\nabla g_b(ru)|^2
 =r^2\bigl(4h_b(u)^2+|\nabla_Sh_b(u)|^2\bigr).
 \tag{1}
\]
The two summands in the gradient are orthogonal. Regularity and
compactness therefore give a uniform \(c>0\) with
\(|\nabla g_b(w)|\geq c|w|\). Upper bounds of the same kind also hold.

Let \(a_b\) be the critical section. A cutoff constant vector field
\(\chi(w)a_b\), equal to \(a_b\) on a common neighborhood of all segments
\([0,a_b]\), has a smooth compactly supported time-one flow \(\phi_b\)
with \(\phi_b(0)=a_b\). Pulling back \(f_b\) by this flow moves its
unique critical point to zero without changing infinity.
The whole flow is a family of diffeomorphisms, so it preserves
nondegeneracy and the number of critical points.

Let \(v_b=f_b(0)\). Subtract
\(t v_b\chi(|w|/A)\), \(0\leq t\leq1\), where \(\chi=1\) on
\([0,1]\), \(\chi=0\) on \([2,\infty)\), and \(|\chi'|\leq C_\chi\).
Take \(A\) beyond all previous compact modifications.
On the only transition annulus the gradient error is at most
\(\sup|v_b|C_\chi/A\), whereas the old gradient is at least \(cA\).
Choosing \(A^2>2\sup|v_b|C_\chi/c\) excludes a new critical point
throughout the homotopy. Elsewhere the correction is constant or zero.
The result has critical value zero.

Write its Taylor expansion as \(f_b=Q_b+R_b\), where
\[
 Q_b(w)=\tfrac12\operatorname{Hess}_0(f_b)(w,w),\quad
 |R_b(w)|\leq K_0|w|^3,\quad |dR_b(w)|\leq K_1|w|^2,
 \quad |\nabla Q_b(w)|\geq c_0|w|.
 \tag{2}
\]
All constants and the Taylor neighborhood are uniform.
Replace \(f_b\) by \(Q_b+(1-t\chi(|w|/a))R_b\).
For \(r=|w|\leq2a\), the error from \(dQ_b\) is bounded by
\[
 (K_1+2C_\chi K_0)r^2.
 \tag{3}
\]
Indeed the differentiated cutoff contributes at most
\(C_\chi K_0r^3/a\), on \(a\leq r\leq2a\).
Take \(4a(K_1+2C_\chi K_0)<c_0\).
Every intermediate gradient is nonzero for \(r>0\) in the modified
ball; outside it the function is unchanged.
Its Hessian at zero is unchanged. We have reached \(f_b=Q_b\)
on one common ball, with no change at infinity.

\begin{lemma}\label{lem:sphere}
Suppose that the preceding local quadratic family has smooth
\(T_{k,l}\) infinity data, \(k,l\geq3\), over a compact smooth base.
Its Morse index is \(k\). If \(E^-\) is its negative eigenspace
bundle in the ambient Euclidean metric, there is a fiber homotopy
equivalence \(S(E^-)\longrightarrow D(h)\) over the base.
The local quadratic forms need not be constant or globally framed.
\end{lemma}
\begin{proof}
For small uniform \(\delta>0\), put \(M_\pm^b=\{f_b\leq\pm\delta\}\).
Choose a large sphere of radius \(R\) in a common homogeneous collar.
On \(M_-^b\cap\{r\leq R\}\) choose a strictly descending field,
radial outward near \(r=R\). This prescription is possible because
\(df_b(\partial_r)=2f_b/r<0\) there. Convex patching with negative
gradient gives a field smooth over all parameters. The compact inner
region contains no critical point. Its bounded function range and a
uniform positive lower bound for \(-df_b(V)\) force exit through
\(r=R\) in uniformly bounded time; exit through \(f_b=-\delta\) is
impossible. The exit is transverse.
Stopped flow, fixing the exterior, is a continuous family deformation
retraction onto the outer part of \(M_-^b\).

Angular projection of that outer part has, over \(h_b(u)<0\), the
radial interval
\[
 \left[\max\left\{R,\sqrt{\delta/(-h_b(u))}\right\},\infty\right).
 \tag{4}
\]
Its continuous lower endpoint gives a contraction of the interval
fibers. The open negative domain is homotopy equivalent to its
closure, uniformly in the parameters, by a small regular boundary
collar. Each negative domain has type \(S^{k-1}\): a path in
\(T_{k,l}\) can be smoothed while preserving regularity, and solving
\(\partial_t h+dh(X)=0\) near its zero sets gives an ambient domain
isotopy. Only this individual fiber type, not a simultaneous choice
of paths, is used. Thus \(M_-^b\simeq S^{k-1}\).

On \(f_b\geq\delta\), normalized descent
\(-\nabla f_b/|\nabla f_b|^2\) has derivative \(-1\) on \(f_b\)
and bounded speed, by (1) and compactness away from the critical
point. Extend it using a cutoff below \(\delta\), away from zero.
It is a complete bounded field. Stopping at level \(\delta\) after
time \(f_b(w)-\delta\), and fixing the lower sublevel, retracts
\(\R^{k+l}\) onto \(M_+^b\). In particular \(M_+^b\) is contractible.

Here is the relative handle comparison, without choosing frames of
the eigenspaces. Split the trivial ambient bundle as \(E^-\oplus E^+\)
by the smooth spectral projections of \(Q\).
In its adapted metrics write \(Q=-|x|_Q^2+|y|_Q^2\).
These metrics and their comparison constants are smooth and uniform.
Choose a smooth \(\nu\) supported in \((0,2\delta)\), with
\(0\leq\nu\leq3/4\), equal to \(3/4\) on
\([\delta/4,7\delta/4]\), and set
\(\mu(s)=\int_s^{2\delta}\nu(u)\,du\).
Then
\[
 -3/4\leq\mu'\leq0,\qquad
 \mu(0)\geq9\delta/8>\delta,\qquad
 \mu=0\ \text{for }s\geq2\delta.
 \tag{5}
\]
In a sufficiently small quadratic neighborhood use
\(\widetilde f=Q-\mu(|x|_Q^2+2|y|_Q^2)\).
Its support is inside \(Q<\delta\), and in orthonormal coordinates
for these metrics its derivatives are
\[
 -2(1+\mu')x,\qquad 2(1-2\mu')y.                 \tag{6}
\]
Thus it has just the original nondegenerate critical point, now
below \(-\delta\), and \(\{\widetilde f\leq\delta\}=M_+\).
Normalized descent across its regular slab retracts \(M_+\) onto
\(\{\widetilde f\leq-\delta\}\), fixing the latter and hence \(M_-\).

The latter set retracts relative to \(M_-\) onto \(M_-\cup D\), where
\(D=\{y=0,\ |x|_Q^2\leq\delta\}\).
At a point outside \(M_-\), shrink \(y\) radially to zero when
\(|x|_Q^2\leq\delta\), and otherwise to norm
\(\sqrt{|x|_Q^2-\delta}\). During the motion \(\widetilde f\)
decreases since its derivative with respect to \(|y|_Q^2\) is
\(1-2\mu'>0\). The formula agrees continuously with fixing \(M_-\).
Membership of \(D\) follows from
\(s+\mu(s)\geq\mu(0)>\delta\) for \(0\leq s\leq\delta\).
All formulas are invariant under the orthogonal transitions of
the adapted bundles.

If the index is \(j\), the pair has the relative homology of
\((D^j,S^{j-1})\). Since \(M_+\) is contractible and
\(M_-\simeq S^{k-1}\), its relative homology is instead \(\Z\) in
degree \(k\), and zero otherwise. Hence \(j=k\), and the boundary
of the local negative disk is a generator of
\(H_{k-1}(M_-^b;\Z)\).
Expand that boundary, within the local negative cone and still in
\(M_-\), to a fixed small Euclidean sphere \(aS(E_b^-)\).
Flow it to \(r=R\) using the preceding descending field and project
angularly. This gives a smooth bundle map into the open negative
domain. On every fiber it takes the generator to a generator.
Both spaces have the type of the simply connected sphere
\(S^{k-1}\), so homological Whitehead makes each fiber map a
homotopy equivalence.

For clarity, the last fiberwise assertion is a family assertion.
For a bundle map with fiber homotopy equivalences over a CW base,
the fiberwise mapping-path replacement has contractible homotopy
fibers. Local bundle trivializations and the relative extension
property over successive cells construct a section and the
fiberwise homotopies to the identities. Equivalently, every
obstruction is a homotopy group of one of these contractible fibers
and vanishes. This is the fiberwise Whitehead argument, not a
selection of independently chosen inverses. It proves the claimed
fiber homotopy equivalence over the whole base.
\end{proof}

\section{Fixing the quadratic model and stabilizing smoothly}

Apply Lemma~\ref{lem:sphere} with \(k=l=p\).
The spherical hypothesis gives \(BJ[E^-]=0\).
Adams's injectivity therefore implies \([E^-]=0\) in \([S^9,BO]\).
The passage between free and based sphere classes is harmless:
the fundamental group action fixes the zero element, so its orbit
contains only zero.

The homotopy fiber of \(\Gr_p(\R^{2p})\to BO(p)\) is
\(V_p(\R^{2p})\), which is \((p-1)\)-connected.
This follows successively from the sphere fibers in the
forgetful maps of Stiefel manifolds.
Also \(BO(p)\to BO(p+1)\) has fiber \(S^p\), so
\(BO(p)\to BO\) is \(p\)-connected.
Both maps induce isomorphisms on \(\pi_9\), since \(p\geq20\).
Thus the actual negative-plane map is null in the finite
Grassmannian, not only after adding an unspecified bundle.

For a symmetric invertible matrix \(A\) of \(Q\), the path
\((1-t)A+tA|A|^{-1}\) remains invertible with the same negative plane.
Sign forms are \(I-2P_E\).
A smooth, endpoint-collared approximation to the Grassmannian
nullhomotopy consequently gives nondegenerate quadratic forms
\(Q_{b,s}\), from \(Q_b\) to the fixed \(q_{p,p}\), with
\[
 |\nabla Q_{b,s}(w)|\geq c_1|w|,\qquad
 |\partial_s Q_{b,s}(w)|\leq K|w|^2.              \tag{7}
\]
Compactness gives uniform constants. No orthogonal lift or
chosen frame of \(E^-\) has been asserted.

Insert \(Q_{b,\lambda s(r)}(w)\), \(0\leq\lambda\leq1\), in the
quadratic ball, with \(s=1\) near zero and \(s=0\) near its outer
boundary. Its differential is the fixed-parameter differential plus
\[
 \lambda\,\partial_sQ_{b,\lambda s(r)}(w)s'(r)\,dr .
 \tag{8}
\]
Choose \(r|s'(r)|\leq\varepsilon<c_1/(2K)\) if \(K>0\);
there is no restriction from this error when \(K=0\).
Such a cutoff is obtained on a logarithmic interval of length \(L\):
its derivative satisfies \(r|s'|\leq\|\sigma'\|_\infty/L\).
The outer radius stays inside the quadratic ball and the inner
radius decreases as \(e^{-L}\). The error in (8) is at most
\(K\varepsilon r<c_1r/2\).
Every intermediate function has exactly its nondegenerate critical
point at zero and the same infinity data. At the endpoint its
local form is the fixed \(q_{p,p}\).

\begin{lemma}[Quantitative stabilization smoothing]\label{lem:smoothing}
Let \(G_0(b,w,z)=g_b(w)+q'(z)\), where \(g_b\) is smooth away
from zero, homogeneous of degree two, and
\(|\nabla g_b(w)|\geq c|w|\), and \(q'\) is fixed nondegenerate
quadratic. Suppose a smooth \(F\) satisfies
\[
 |F-G_0|\leq M_0,\qquad |dF-dG_0|\leq M_1 .
 \]
Over any parameter set for which the bounds below are uniform,
one can modify \(F\) outside an arbitrarily large common ball,
through smooth functions with no new critical points, to a smooth
homogeneous function at infinity. All original critical
neighborhoods are unchanged. The new homogeneous data are
regularly \(C^1\)-homotopic to \(G_0\).
\end{lemma}
\begin{proof}
Put \(R=|(w,z)|\), \(a=|w|\), \(s=a/R\).
Block orthogonality gives
\(|\nabla G_0|\geq c_*R\) for some \(c_*>0\).
Choose uniform bounds
\(|g_b(w)|\leq H_0a^2\), \(|dg_b(w)|\leq H_1a\).
Let \(\rho_\eta(s)=\rho(s/\eta)\), with \(\rho=0\) on
\([0,1]\), \(\rho=1\) on \([2,\infty)\), \(0\leq\rho\leq1\),
and \(|\rho'|\leq K_\rho\). Set
\[
 G_\eta=\rho_\eta(a/R)g_b(w)+q'(z).
 \]
The possibly nonsmooth summand vanishes on an open cone about \(w=0\),
so \(G_\eta\) is smooth away from the total origin.
Where the cutoff changes,
\[
 |ds|^2=\frac{|z|^2}{R^4}\leq R^{-2}.
 \tag{9}
\]
On the support of \(G_\eta-G_0\), \(a\leq2\eta R\), giving
\[
 |G_\eta-G_0|\leq D_0\eta^2R^2,\quad
 |dG_\eta-dG_0|\leq D_1\eta R,\quad
 D_0=4H_0,\quad D_1=2H_1+4K_\rho H_0 .
 \tag{10}
\]
These bounds extend across \(w=0\) by \(C^1\) continuity.
For \(D_1\eta<c_*/2\), straight interpolation between \(G_0\)
and \(G_\eta\) has no nonzero critical point, hence regular angular
zero level by (1).

Choose \(\alpha_A(R)=\alpha(R/A)\), zero for \(R\leq A\), one
for \(R\geq2A\), with \(0\leq\alpha\leq1\),
\(|\alpha'|\leq K_\alpha\). For all \(0\leq\tau\leq1\) use
\[
 F_\tau=(1-\tau\alpha_A)F+\tau\alpha_AG_\eta .
 \tag{11}
\]
On \(A\leq R\leq2A\), differentiating the cutoff as well as the
two functions bounds the gradient error from \(G_0\), divided by \(R\),
by
\[
 \frac{M_1}{A}+D_1\eta+
 \frac{K_\alpha M_0}{A^2}+2K_\alpha D_0\eta^2 .
 \tag{12}
\]
For example, the last term comes from
\((K_\alpha/A)D_0\eta^2R^2/R\), using \(R/A\leq2\);
it must not be discarded. The value-error term is of order \(A^{-2}\),
not merely an unspecified small perturbation.
Choose \(\eta>0\), at most \(1/4\), so that
\[
 \frac{D_1\eta}{c_*}\leq\frac18,\qquad
 \frac{2K_\alpha D_0\eta^2}{c_*}\leq\frac18,
 \]
and then choose \(A\), as large as desired, so that
\[
 \frac{M_1}{c_*A}\leq\frac18,\qquad
 \frac{K_\alpha M_0}{c_*A^2}\leq\frac18 .
 \tag{13}
\]
If a numerator vanishes its inequality imposes no condition.
Then (12) is at most \(c_*/2\).
Outside \(2A\) the cutoff-derivative terms vanish and the same
lower bound applies throughout (11). Every original critical point
lies in \(R\leq M_1/c_*\); enlarge \(A\) past a neighborhood of this
ball. Thus the homotopy creates none and preserves all of them,
including their Hessians. At \(\tau=1\) it equals \(G_\eta\)
outside \(2A\), and it remains smooth at the origin because it is
unchanged there.

An intermediate \(F_\tau\) need not be exactly homogeneous outside
any ball. Its homogeneous reference is
\(G_\tau=(1-\tau)G_0+\tau G_\eta\), and outside \(2A\)
\(F_\tau-G_\tau=(1-\tau)(F-G_0)\), with bounded value and derivative.
Thus the associated explicit regular homogeneous homotopy tracks
the stabilized tube class. Exact smooth homogeneity, needed for the
zero-band theorem, is restored at \(\tau=1\); that theorem is not
applied at intermediate times.
\end{proof}

Add exactly ten positive squares to the normalized \(f_b\).
The direct sum has just its nondegenerate critical point at zero.
Its difference from \(g_b(w)+|z|^2\), including its vertical derivative,
is uniformly bounded: the difference is supported in a fixed
\(w\)-ball, although not in a fixed \((w,z)\)-ball.
Lemma~\ref{lem:smoothing} applies and leaves the fixed local form intact.
We have obtained
\[
 k=p,\qquad l=p+10,\qquad l>k+9.                    \tag{14}
\]
The smoothed angular data have type \(T_{k,l}\). Indeed direct
sums of a \(C^1\) regular path from the original datum to
\(q_{p,p}\) remain regular, by block gradients, and the smoothing
has just been joined to that direct sum by a regular homotopy.
This argument preserves the actual varying stable infinity class;
it has not replaced it by fixed quadratic infinity data.

\section{From the accepted product complement to unframed cores}

Apply the accepted zero-band result to this normalized family.
It gives smooth embeddings \(j_b:T^0\to\Int T_b^1\) and an
incoming-relative product identification of
\(T_b^1\setminus\Int j_b(T^0)\) over the whole \(S^9\).
The numerical hypothesis is particularly safe here:
\[
 \dim\Sigma=k+l-2\geq48>35,\qquad
 \min(k,l)-2\geq18>9 .
 \tag{15}
\]
These are only the arithmetic applications of the already
reviewed connectivity argument.

Move the actual negative boundary through its small regular angular
collar to the truncated one. Then, in the \emph{family} product
coordinate on the attached cobordism, shorten \([0,1]\) to
\([0,t]\), \(1\geq t\geq0\). At \(t=0\) its boundary is
\(\partial j_b(T^0)\), an embedded hypersurface.
The fixed incoming collars make this an isotopy of domains smooth
also at \(t=0\): extend the product coordinate slightly into
\(j_b(T^0)\). This uses the accepted family trivialization, not
the fact that each cobordism separately is a product.

Fix \(T^0=S^{k-1}\times D^l\) and set \(e_b(u)=j_b(u,0)\).
Choose the metric normal bundle \(\nu\) of this family of cores.
Compactness gives a common tubular chart. Shrink the disk factors
of \(j_b\) to radius \(r_0>0\), small enough that their images
lie in this chart, and express them there as
\[
 j_b(u,z)=(\psi_b(u,z),\xi_b(u,z)),\quad
 \psi_b(u,0)=u,\quad \xi_b(u,0)=0,\quad |z|\leq r_0 .
 \]
Here \(\xi_b(u,z)\) belongs to the normal fiber over
\(\psi_b(u,z)\); no identification of different normal fibers has
been omitted. For \(0<t\leq1\) use
\[
 L_t(b,u,z)=\bigl(\psi_b(u,tz),\,t^{-1}\xi_b(u,tz)\bigr).
 \tag{16}
\]
It is a composition of source disk dilation, the original embedding,
and fiberwise normal-bundle dilation, so it is an embedding.
Uniformly \(|\xi_b(u,tz)|\leq Ct|z|\); take \(Cr_0\) smaller
than the common normal-chart radius.

In any local trivialization of \(\nu\), the normal coordinates satisfy
the exact division formula
\[
 t^{-1}\xi_b(u,tz)
   =\int_0^1D_z\xi_b(u,vtz)\,z\,dv .                \tag{17}
\]
Thus (16) extends smoothly at \(t=0\) to
\((u,A_b(u)z)\), \(A_b(u)=D_z\xi_b(u,0)\).
On chart overlaps the smooth limits agree with (16) for \(t>0\)
and hence agree at zero. The normal derivative \(A_b(u)\) is an
isomorphism: the derivative of the full codimension-zero embedding
at the core induces an isomorphism on the quotient by its tangent
space. Its limit is therefore also an embedding.

The resulting normal ellipsoid is
\(\{v:\langle(AA^*)^{-1}v,v\rangle\leq r_0^2\}\).
Interpolate the positive definite operator \((AA^*)^{-1}\) with
the identity. Compactness gives a positive common lower eigenvalue
throughout this interpolation, so a further small choice of \(r_0\)
keeps every ellipsoid in the same chart. Its boundary is a smooth
embedded sphere bundle at every time. We have deformed the domains
to round normal neighborhoods of \(e_b\).
This deformation forgets the parametrization of the disk fibers.
It neither chooses nor asks the ensuing core homotopy to preserve
the frame \(A_b\).

\begin{lemma}[Relative parameter-uniform general position]
\label{lem:relative}
Let \(P=B\times[0,1]\), with \(B\) a closed compact smooth
manifold, \(\dim P=q\), and let \(M\) be a closed compact
\(m\)-manifold. Suppose a smooth
\(H:P\times M\to Y^N\) is constant in time and has fiber embeddings
on collars of the two time endpoints, and \(N>q+2m\).
Then it has an arbitrarily small smooth perturbation, fixed on
smaller endpoint collars, which is an embedding in every fiber.
\end{lemma}
\begin{proof}
Double the time interval by concatenating \(H\) with
its reverse. This gives a smooth map
\(\overline H:\overline P\times M\to Y\) on
\(\overline P=B\times S^1\), since the collars are constant.
We can work throughout on closed compact source manifolds.

Write \(X=\overline P\times M\). In \(J^1(X,Y)\), the stratum
where the derivative restricted to the vertical tangent \(TM\)
has rank \(r<m\) has codimension
\((N-r)(m-r)\), at least \(N-m+1\).
The other derivative columns are unrestricted. Since
\[
 N-m+1>q+m=\dim X ,
 \tag{18}
\]
transversality to these finitely many strata makes their inverse
images empty.

For double points use the ordinary two-point zero-jet space
\(J_2^0(X,Y)\), with its distinct-source-point condition.
The \emph{bad} submanifold consists of pairs
\[
 ((p,u),(p,v);y,y),\qquad u\ne v .
 \]
Equality of the two parameter coordinates has codimension \(q\);
equality of their target values has codimension \(N\).
These are independent conditions, so its codimension is \(q+N\).
Its inverse image under the two-point jet of a transverse map is
empty because
\[
 2(q+m)-(q+N)=q+2m-N<0 .                          \tag{19}
\]
This forbids only \emph{same-parameter} double points.
Imposing ordinary double-point avoidance without the equal-parameter
condition would use the wrong, stronger dimension inequality.

The multijet theorem in the stated input, applied to all these
strata, gives a dense residual set of maps \(G\) satisfying both
avoidances. Choose \(G\) arbitrarily close to \(\overline H\) in
the smooth topology. On a central closed subinterval of the
original time interval we keep \(G\).
Take a smooth parameter cutoff \(\beta\), equal to one there and
zero on smaller endpoint collars, whose transition regions lie
inside the original embedding collars.
For \(G\) sufficiently close to \(\overline H\), set on the original
cylinder
\[
 H'(p,u)=\exp_{H(p,u)}
       \bigl(\beta(p)\log_{H(p,u)}G(p,u)\bigr).
 \tag{20}
\]
The logarithm is defined on a uniform neighborhood of the compact
image. Where \(\beta=1\), (18)--(19) give injective immersions.
Where \(\beta\) changes, the original fiber maps are embeddings,
and (20) is uniformly \(C^1\)-close to them. Embeddings of a compact
source are \(C^1\)-open, uniformly over a compact family:
the derivative has a common lower singular-value bound; finitely
many source charts give uniform local injectivity; away from a
small diagonal neighborhood compactness gives positive separation.
These same bounds hold for sufficiently close maps.
Thus (20) is a fiber embedding also on the transition collars.
It is exactly \(H\) on the prescribed smaller collars.
An injective immersion of compact \(M\) is an embedding.
This proves the relative assertion; no relative conclusion was
silently imported from the non-relative corollary.
\end{proof}

For our cores \(M=S^{k-1}\) and \(Y=S^{k+l-1}\).
As maps on \(S^9\times S^{k-1}\), both \(e_b(u)\) and the standard
core family are nullhomotopic: that finite CW source has dimension
\(k+8<k+l-1\), while the target sphere is
\((k+l-2)\)-connected. Cell-by-cell extension, or cellular
approximation, gives the nullhomotopies.
Concatenate them and smooth relative to endpoint collars to obtain
an initial homotopy \(H\).
Here \(q=10\), \(m=k-1\), \(N=k+l-1\), and
\[
 N-(q+2m)=l-k-9\geq1,\qquad
 (N-m+1)-(q+m)=l-k-8\geq2 .                       \tag{21}
\]
Lemma~\ref{lem:relative} supplies the needed homotopy through
embeddings, fixed at both ends.

The metric normal bundle over the compact
\(S^9\times[0,1]\times S^{k-1}\) has a uniform tubular radius.
For completeness, near its zero section the normal exponential
has invertible derivative; compactness gives uniform local charts.
If arbitrarily small radii still allowed two normal points in one
fiber with the same image, take a convergent sequence of their base
points. The limiting core images agree, so the base points agree
by injectivity of that fiber core. They then lie in one of the
uniform local charts, contradicting its injectivity.
Thus one positive radius works for the entire family.
Shrink our initial round tube to this radius if necessary, and
follow these normal disk \emph{bundles} along the homotopy.
At the end the core is the fixed standard core, so its round
neighborhood is fixed. Increasing its angular radius to \(\pi/4\)
identifies it by a domain isotopy with \(D(q_{k,l})\).
No normal framing has been extended through the homotopy.

\section{Functions, infinity classes, and conclusion}

A smooth compact family of cooriented embedded domain boundaries
has a uniform signed normal collar. A smooth profile, equal to
the signed coordinate near zero and to fixed negative and positive
constants on the two outer portions, gives a smooth family of
defining functions on \(Y\).
For a fixed domain their space is convex: signs in the two interiors
and the positive outward derivative on the boundary persist under
convex combination. This also permits a prescribed initial or final
defining function, by an extra endpoint-collared convex interpolation.
Consequently all the preceding domain homotopies lift to regular
smooth angular function homotopies, with the original infinity
datum at the start and \(q_{k,l}|_Y\) at the end.
They prove nullity in \(T_{k,l}\), not merely fiberwise diffeomorphism
of the negative domains.

These last homotopies need only be homotopies of infinity data.
Nevertheless they can, if desired, also be realized by smooth
one-critical families without changing the critical neighborhood.
For a smooth regular homogeneous homotopy \(G_{b,t}\), compactness
gives \(|d_wG_{b,t}|\geq c'|w|\) and
\(|\partial_tG_{b,t}|\leq K'|w|^2\).
Starting in a homogeneous outer region, set
\[
 F_{b,\lambda}(w)=G_{b,\lambda s(|w|)}(w),
 \tag{22}
\]
with \(s=0\) near its inner radius, \(s=1\) outside a much larger
radius, and \(r|s'(r)|<c'/(2K')\) when \(K'>0\).
Keep the original function inside. The same calculation (8)
gives nonzero gradient on the new annulus, uniformly in \(b,\lambda\).
The critical point, its nondegeneracy, and its neighborhood are
unchanged. Outside the new common ball the data are \(G_{b,\lambda}\).
Thus even this stronger realization does not introduce an
uncontrolled critical point or replace the stable infinity class.

Finally add ten negative squares. With the fixed ordering of
negative coordinates before positive coordinates this reaches
the diagonal stage \(T_{p+10,p+10}\).
The composite addition is the original datum stabilized by ten
copies of \(q_{1,1}\), up to a fixed coordinate permutation.
That permutation is a homeomorphism of function spaces carrying
constant data to constant data; no claim that it is isotopic to
the identity is needed for nullity.
Adding a fixed quadratic summand preserves regularity of the
homogeneous homotopy by block gradients. Angular smoothing, or
Lemma~\ref{lem:smoothing} if smooth total functions are wanted,
preserves its class.
The nullhomotopy at \(T_{k,l}\) therefore proves stable nullity
of the original family. This proves Theorem~\ref{thm:main}.

\paragraph{What improved over the prior result.}
The accepted zero-band theorem began with a fixed local quadratic
model and ended with triviality of a cobordism family.
Here the original varying negative eigenspace bundle is compared
over the whole base with infinity, \(J_8\) makes that \emph{finite}
local model null, all changes of functions have quantitative
gradient margins, and the product cobordism is removed followed by
a genuinely relative homotopy of unframed cores.
These close the remaining logical dependencies of
\texttt{MS-one-critical}; they do not prove the nonzero-kernel
proposition or the manuscript's main theorem.

\paragraph{Exact-certificate boundary.}
\texttt{verify.py} checks the local cutoff and radial chain rules,
the handle derivatives and their margins, the angular ratio
differential (9), all terms in (10)--(13), the block-gradient bound,
a polynomial Taylor-jet replay of (16), and the universal dimension
polynomials for \(k=20+u\), \(l=k+10+v\), \(u,v\geq0\).
One bounded Jacobian operation encloses
\(1-e_1-e_2-e_3-e_4\) on \([0,1/8]^4\), independently replayed
over the rationals. The four variables are exactly the four
normalized error budgets in (13), not arbitrary proxies for the
gradient. Its range is \([1/2,1]\), with both endpoints attained.
This is a certificate for the arithmetic after the analytic
estimates, not a theorem about an arbitrary function supplied to
a computer.
Neither it nor the polynomial normal-coordinate example verifies
the smooth bundle arguments, ODE constructions, relative
transversality, the external \(J\)-theorem, or the accepted
h-cobordism theorem. Those are the prose proof and explicit inputs.
No independent approval of this new proof is claimed.

\begin{thebibliography}{9}
\bibitem{Adams}
J.~F. Adams, \emph{On the groups \(J(X)\). IV},
Topology 5 (1966), 21--71.
\href{https://doi.org/10.1016/0040-9383(66)90004-8}
{doi:10.1016/0040-9383(66)90004-8}.
Original retained from
\url{https://math.univ-cotedazur.fr/u/cazanave/Gdt/ImJ/J-IV.pdf}.
\bibitem{BPT}
M. Borodzik, M. Powell and P. Teichner,
\emph{Link concordance implies link homotopy},
\href{https://arxiv.org/abs/2503.04348v1}{arXiv:2503.04348v1},
Theorem 2.3.4 and Corollary 2.3.5(iii).
\bibitem{Hatcher}
A. Hatcher, \emph{Algebraic Topology}, Cambridge University Press,
2002, ISBN 9780521795401; corrected author electronic text,
Corollary 4.33. Retained in the preceding zero-band audit.
\end{thebibliography}
\end{document}
